%%%%%%%%%%%%%%%%%%%%%%%%%%%%%%%%%%%%%%%%%%%%%%%%%%%%%%%
% % Lines starting with % are comments, which are ignored.
% % This is a handy way of indicating the date and version of
% % your document, to wit:
% %
% % LaTeX sample file
% % Modified March, 2002
% %
%%%%%%%%%%%%%%%%%%%%%%%%%%%%%%%%%%%%%%%%%%%%%%%%%%%%%%%
\DocumentMetadata{
  % PDF standards: ONE key whose value is a comma-separated list --
  % UA-2 (accessibility) plus A-4f (archival; the archival companion of
  % UA-2). Do NOT use separate pdfstandard=... keys: with ordinary
  % key-value parsing only the last one takes effect.
  pdfstandard={UA-2,A-4f},
  lang=en-US,
  tagging=on,
  tagging-setup={
    math/setup=mathml-SE
  }
}

\documentclass[12pt]{article}
% homework-macros.tex -- shared package setup for homework-template.tex
% (article class).
%
% NOTE ON PATHS: every template in this repository is stored one level up
% from where it is compiled (see README.md). The \input path below is
% resolved against the folder the document is compiled in, which sits
% one directory above this file.

% amsmath + amssymb + amsthm: full math support and proof environments
% (\newtheorem, \begin{proof}, ...).
\usepackage{amsmath, amssymb, amsthm}

% unicode-math: math typeset with OpenType math fonts (needed under
% XeLaTeX/LuaLaTeX so math matches the text fonts).
\usepackage{unicode-math}

% hyperref: clickable links (the template uses \href for the textbook
% problem sources); loaded after the math packages so it can patch them.
% Uses hyperref's default link colors.
\usepackage[colorlinks=true]{hyperref}

% Shared course math shorthands (\RR, \NN, \ZZ, \QQ, \grad, \curl, \tr,
% \der, \pder, ...) from the repository root.
\newcommand{\RR}{\mathbb{R}}
\newcommand{\NN}{\mathbb{N}}
\newcommand{\ZZ}{\mathbb{Z}}
\newcommand{\QQ}{\mathbb{Q}}
\newcommand{\vv}{\vec{v}}
\newcommand{\ww}{\vec{w}}
\newcommand{\xx}{\vec{x}}
\newcommand{\yy}{\vec{y}}
\newcommand{\zz}{\vec{z}}
\newcommand{\uu}{\vec{u}}
\newcommand{\basis}[1]{\vec{e_{#1}}}
\newcommand{\id}{\mathrm{id}}
\newcommand{\CC}{\mathbb{C}}
\renewcommand{\div}{\operatorname{div}}
\newcommand{\grad}{\operatorname{grad}}
\newcommand{\curl}{\operatorname{curl}}
\newcommand{\tr}{\operatorname{tr}}
\newcommand{\der}[2]{\dfrac{d #1}{d #2}}
\newcommand{\pder}[3][]{\dfrac{\partial^{#1} #2}{\partial #3}}


% % Title and author(s)
%%%%%%%%%%%%%%%%%%%%%%%%%%%%%%%%%%%%%%%%%%%%%%%%%%%%%%%
\title{Homework 5 - 09/22 Content}
\author{MATH 1271, Garrett Credi}
\date{09/22/2026}
%%%%%%%%%%%%%%%%%%%%%%%%%%%%%%%%%%%%%%%%%%%%%%%%%%%%%%%
% %
% % The next command allows your in import encapsulated
% % postscript files, .epsf or .eps files, which
% % contain vector graphic image data.
% %
%%%%%%%%%%%%%%%%%%%%%%%%%%%%%%%%%%%%%%%%%%%%%%%%%%%%%%%
% %   The first thanks indicates your affiliation
% %
% %  Just the name here.
% %
% % Your mailing address goes at the end.
% %
% % \thanks is also how you indicate grant support
% %
%%%%%%%%%%%%%%%%%%%%%%%%%%%%%%%%%%%%%%%%%%%%%%%%%%%%%%%

%%%%%%%%%%%%%%%%%%%%%%%%%%%%%%%%%%%%%%%%%%%%%%%%%%%%%%%
% $ The next command ensures the margins are 1 inch on
% % all sides.
% %
%%%%%%%%%%%%%%%%%%%%%%%%%%%%%%%%%%%%%%%%%%%%%%%%%%%%%%%
\begin{document}
\maketitle

Problems listed as AC are found in the \href{https://activecalculus.org/single1e/frontmatter.html}{Active Calculus} 
textbook at the end of the section listed.

Problems listed as APEX are found in the \href{https://opentext.uleth.ca/apex-calculus/apex-calculus.html}{APEX Calculus Textbook}
again at the end of the listed section.

Problems listed as PAULS-PRACT are found in \href{https://tutorial.math.lamar.edu/Problems/CalcI/CalcI.aspx}{Paul's Online Math Notes - Review Problems}
while problems listed as PAULS-ASSIGN are found in \href{https://tutorial.math.lamar.edu/ProblemsNS/CalcI/CalcI.aspx}{Paul's Online Math Notes - Assignment Problems}
and are indexed via [Topic]-[Section]-[Problem Number]

\section{Type 1 - Algebraically Computing the Derivative}
Choose 4 of the following problems to solve.
\begin{enumerate}
  \item AC 1.3.3
  \item AC 1.3.8
  \item AC 1.3.9
  \item APEX 2.1.7-14 (each of these counts as one problem)
  \item PAULS-ASSIGN Derivatives.The Definition of the Derivative.1-24 (each of these counts as one problem)
\end{enumerate}

\section{Type 2 - Graphically Interpreting the Derivative}
Choose 4 of the following problems to solve.

\begin{enumerate}
  \item AC 1.3.2
  \item AC 1.3.4
  \item AC 1.3.5
  \item APEX 2.1.29-32 (each of these counts as one problem)
  \item APEX 2.1.33-34 (each of these counts as one problem)
\end{enumerate}

\section{Type 3 - Sketching Graphs}
Choose 2 of the following problems to solve.

\begin{enumerate}
  \item PAULS-ASSIGN.Derivatives.Interpretations of the Derivative.4-6 (each of these counts as one problem)
  \item PAULS-ASSIGN.Derivatives.Interpretations of the Derivative.7-9 (each of these counts as one problem) 
\end{enumerate}

\section{Type 4 - Algebraically Interpreting the Derivative}
Choose 3 of the following problems to solve.

\begin{enumerate}
  \item PAULS-ASSIGN.Derivatives.Interpretations of the Derivative.10,12,13,14,16,17,18,21 (each of these counts as one problem)
\end{enumerate}

\section{Type 5 - Tangent Lines}
Choose 2 of the following problems to solve.

\begin{enumerate}
  \item AC 1.3.6
  \item AC 1.3.7
  \item APEX 2.1.15-22 (each of these counts as one problem, ignore the bit about the ``normal line")
  \item PAULS-ASSIGN.Derivatives.Interpretations of the Derivative.11,15,19 (each of these counts as one problem)
\end{enumerate}

\section{Challenge Problem}

From our discussion on the squeeze theorem we know that, for $f(x) = \sin(x)$:
\begin{equation}
  f'(0) = \lim_{x \rightarrow 0} \dfrac{\sin(x)}{x} = 1
\end{equation}

because, when $x \approx 0$ we had that $\cos(x) \leq \dfrac{\sin(x)}{x} \leq 1$.
Since this relies of $x \approx 0$, a similar strategy of using that inequality 
and the squeeze theorem will \textit{not} let us compute $f'(a)$ for $a \neq 0$.

However, you \textit{do} have the tools necessary to compute $f'(a)$ for \textit{every} $a \neq 0$ - 
this is your goal for this challenge problem: to compute the derivative of $\sin(x)$ at every $x = a \neq 0$.

If you get stuck, I have placed increasingly revealing hints on the subsequent pages. 
I highly recommend avoiding the hints as much as possible, since I believe
this is one of the most attainable results for you all to feel the feeling of ``discovery."

\newpage
Hint 1: You might find the 
\begin{equation}
  f'(a) = \lim_{h \rightarrow 0} \dfrac{f(a + h) - f(a)}{h}
\end{equation}
form of the derivative more useful here.
\newpage
Hint 2: Since we know $\lim_{h \rightarrow 0} \dfrac{\sin(h)}{h}$, it might be a good idea to try and reduce our hard limit
to the easier case of the aforementioned limit plus maybe some other terms we have to deal with.
\newpage
Hint 3: $\sin(\alpha + \beta) = \sin(\alpha) \cos(\beta) + \cos(\alpha) \sin(\beta)$.
\newpage
Hint 4: Your computation of limits on worksheet 4 might be useful here.
\newpage
Hint 5: 
\begin{equation}
  \dfrac{\cos(h) - 1}{h} = -\dfrac{1 - \cos(h)}{h^2} \cdot h
\end{equation}
% Most people don't use these, so they are "commented out"
% by starting the lines with a "%"
%\begin{keywords}
%   \LaTeX, typesetting
%\end{keywords}

%\begin{AMS}
%   50C60, 18C25
%\end{AMS}

%%%%%%%%%%%%%%%%%%%%%%
% % Here is the start of the Text
%%%%%%%%%%%%%%%%%%%%%%

\end{document}
